% Data-centric AI comparison report (raw vs structured F1 telemetry)
% Compile: pdflatex comparison_report.tex

\documentclass[11pt,a4paper]{article}
\usepackage[margin=1in]{geometry}
\usepackage[T1]{fontenc}
\usepackage[utf8]{inputenc}
\usepackage{lmodern}
\usepackage{microtype}
\usepackage{booktabs}
\usepackage{enumitem}
\usepackage[hidelinks]{hyperref}

\title{Data-Centric AI: Comparing Lap Analysis from Raw versus Structured F1 Telemetry\\[0.3em]
\large Piastri $\cdot$ Japan $\cdot$ 2026 (single-lap samples)}
\author{}\date{}

\begin{document}
\maketitle

\begin{abstract}
This note compares two analyses of materially the same single-lap telemetry: one confined to raw merged exports, one augmented with project-defined engineered columns. The downstream analytic protocol (persona, refusal to invent session facts) is held similar; the controlled variable is \emph{tabular representation}---in the spirit of data-centric AI, not ``model ingenuity.''
\end{abstract}

\section{Purpose and framing}

This report compares two analyses of what is materially the same single-lap telemetry---one confined to the raw export \path{piastri_2026_japan_fastest_lap_telemetry.csv} (\textbf{raw}), one using \path{piastri_2026_japan_structured.csv}, which adds heuristic engineering columns: \texttt{Corner\_Type}, \texttt{Driving\_State}, and distance-based deciles labelled \texttt{Mini\_Sector}. The downstream analytic tool (prompt and LLM persona) was intentionally similar: evidence-grounded narration, refusal to infer official lap or sector timing, tyre compound, weather, or other facts absent from the CSVs. The \textbf{controlled variable}, in a \textbf{data-centric AI} framing, is therefore \textbf{feature and table design}---not a change of foundation-model family.

\section{Dataset differences beyond column count}

\paragraph{Raw telemetry.}
The raw file exposes \texttt{Speed}, throttle, brake, cumulative \texttt{Distance}, merge metadata (\texttt{Source}), and coordinates. Meaningful segmentation requires explicit analytic choices---for example quartiles by row index (Q1--Q4 as a heuristic) and slow-speed clusters cited by metre along \texttt{Distance}. The narrative can \emph{withhold} named corners because circuit geography is not keyed as text columns in this export.

\paragraph{Structured telemetry.}
The structured file \textbf{reuses identical sensor traces} yet relocates cognition from free-form slicing to preprocessing: categorical summaries of longitudinal intent (\texttt{Driving\_State} counts), coarse corner-style buckets (\texttt{Corner\_Type}), and ten distance bins (\texttt{Mini\_Sector}), enabling tabular summaries per mini-sector instead of hand-built cuts.

\section{Aligned empirical anchors}

Across both summaries, headline statistics \textbf{align}: \textbf{Speed} spans roughly $73$--$317$~km/h (FastF1-style magnitude); approximately \textbf{19\%} of samples have \texttt{Brake=True}; mean \texttt{Throttle} $\approx 65.8$; elapsed sample span near \textbf{93\,s}; the \textbf{DRS} column reads as uniformly inactive in this excerpt. Alignment on these anchors supports consistent empirical grounding despite different narrative affordances.

\section{Outcomes: affordances versus epistemic risk}

\subsection{What raw emphasizes}

Raw-grounded prose excels at \textbf{calling out limits}: normalized \texttt{RelativeDistance} not reaching $1.0$; potential lap-boundary artefacts (e.g., high throttle simultaneously with braking at the earliest rows); ambiguity between interpolation and genuine coasting. These caveats are intellectually conservative and reduce over-claiming about driver technique.

\subsection{What structured emphasizes}

Structured analysis offers \textbf{immediate segment readability}: histogram-style views tying \texttt{Driving\_State} and heuristic corner labels to each \texttt{Mini\_Sector}, enabling shorthand such as dominance of transitional phases or comparatively high mean speed in bands like \texttt{MS\_8}--\texttt{MS\_9} without manual binning repetition. Fluency rises because labels collapse repeated statistical work into categorical columns---but each label inherits upstream threshold choices; readers may mistakenly equate heuristic tags with circuit truth or timing authority unless caveats persist.

\section{Data-centric takeaway}

Improvement attributable to switching inputs stems chiefly from \textbf{upstream representation}: engineered columns compress repeatable rules into shorthand that the same analytic scaffold can cite directly. That productivity gain is \textbf{not} immunity from telemetry limits (boolean brake, sparse merge), nor does it embed lateral tyre physics absent from columns.

\section{Limitations statement}

Neither table supports rival deltas, stint strategy, or optimised compound selection without additional data. Robust reporting should cross-check aggregates with a small scripted audit (ranges, percentages, grouped counts).

\paragraph{Disclaimer.}
Structured labels (\texttt{Corner\_Type}, etc.) remain \emph{project-defined}; distance mini-sectors are not FIA sector timers. Narrative flair must remain accountable to CSV evidence.

\end{document}
